% ==============================================================================
% Informe de Laboratorio - Práctica N.º 1
% Osciloscopio, generador de funciones y voltaje senoidal
% LIE-311 Laboratorio de Circuitos Eléctricos II
% ==============================================================================
%
% CÓMO ENCONTRAR LO QUE FALTA POR LLENAR
% --------------------------------------
% Todo lo que depende de tus mediciones o capturas está marcado con un
% comentario que empieza con:
%
%       % >>> TODO (LLENAR): ...
%
% Busca "TODO" en tu editor para recorrerlos todos. En el PDF, los datos
% faltantes aparecen como líneas en blanco (\vacio) y las capturas del
% osciloscopio como cuadrículas con el texto "Captura pendiente".
%
% Los valores TEÓRICOS (que salen solo de los datos de las tablas 1 y 2 y de las
% señales del montaje 3) ya están calculados. Revísalos y deja tu evidencia de
% cálculo, tal como pide el manual.
% ==============================================================================

\documentclass{article}

% ====== Paquetes básicos ======
\usepackage[utf8]{inputenc}    % Codificación UTF-8
\usepackage[T1]{fontenc}       % Acentos correctos
\usepackage[spanish]{babel}    % Traducción al español
\usepackage{csquotes}          % Recomendado para biblatex con babel
\usepackage{hyperref}          % Hipervínculos
\usepackage{amsmath,amssymb}   % Matemáticas
\usepackage{graphicx}          % Imágenes
\usepackage{caption}           % Mejor control de captions
\usepackage{geometry}          % Márgenes
\usepackage{fancyhdr}          % Encabezados y pies
\usepackage{parskip}           % Mejor separación entre párrafos

% ====== Referencias ======
\usepackage[
    backend=biber,
    style=numeric-comp,     % cambie a APA de ser necesario
    sorting=none
]{biblatex} 
\bibliography{references}

% agrege `\printbibliography{}` para las referencias
% cite con `\cite{}`
% agregue las referencias a un `references.bib`

% ====== Configuración ======
\geometry{margin=2.5cm}
\pagestyle{fancy}
\fancyhf{}
\rhead{\thepage}
\setlength{\headheight}{13pt}
\lhead{\leftmark}

% ====== Comando personalizado para imágenes ======
% Uso: \imagen{ruta}{Texto del pie de figura}
\newcommand{\imagen}[2]{
    \begin{figure}[h!]
        \centering
        \includegraphics[width=200px]{#1}
        \caption{#2}
    \end{figure}
}

% ====== Datos del artículo ======
\title{Portada}
\author{Lester Iván Hernández Hernandez}


% ====== Ayudas locales para este informe ======

% Espacio en blanco para un dato por llenar. Uso: \vacio  o  \vacio[3cm]
\newcommand{\vacio}[1][2cm]{\underline{\hspace{#1}}}

% Recuadro de captura pendiente (pantalla de 8 x 6 divisiones).
% Uso: \capturapendiente{Pie de figura}{Texto dentro del recuadro}{fig:etiqueta}
% Cuando tengas la captura real, reemplaza la llamada por:
%   \imagen{./Pictures/Practica1/nombre.png}{Pie de figura}{475px}
% (y agrega \label{fig:etiqueta} dentro del pie si necesitas referenciarla)
\newcommand{\capturapendiente}[3]{%
    \begin{figure}[ht!]
        \centering
        \begin{tikzpicture}[scale=0.85]
            \draw[step=1cm, gray!60, very thin] (0,0) grid (8,6);
            \draw[thick] (0,0) rectangle (8,6);
            \node[align=center, text=gray] at (4,3) {\textit{Captura pendiente}\\[2pt] \small #2};
        \end{tikzpicture}
        \caption{#1}\label{#3}
    \end{figure}%
}

\begin{document}
    \section{Introducción}

    El osciloscopio y el generador de funciones son dos instrumentos fundamentales en el estudio de las señales periódicas. El primero permite visualizar y analizar señales en el dominio del tiempo, mientras que el segundo permite producir señales con forma, frecuencia y amplitud ajustables. Entre las formas de onda disponibles, la señal senoidal es una de las más utilizadas en electrónica e ingeniería eléctrica, por lo que conocer sus parámetros característicos es indispensable.

    Aunque la teoría define con precisión el valor pico, el valor pico a pico, el período y el valor eficaz de una señal, es necesario comprobar experimentalmente que las lecturas de los instrumentos coinciden con lo esperado y entender qué información entrega cada uno de ellos.

    Por ello, esta práctica tiene como objetivo familiarizarse con el uso del osciloscopio y del generador de funciones, medir los parámetros de señales senoidales y triangulares, y comparar los valores obtenidos con los calculados teóricamente mediante el error relativo porcentual.

    \section{Objetivos}

    \begin{enumerate}
        \item Conocer el generador de funciones y su funcionamiento básico.
        \item Conocer el osciloscopio, su funcionamiento básico y la información que se puede obtener al realizar mediciones con este en los circuitos eléctricos.
        \item Identificar el valor pico, el valor pico a pico y el valor eficaz de una señal senoidal de voltaje.
        \item Comprobar el valor eficaz en una señal senoidal.
    \end{enumerate}

    \section{Materiales y Equipo}

    \begin{enumerate}
        \item Generador de funciones.
        \item Osciloscopio.
        \item Multímetro digital.
        \item Resistencias.
        \item Protoboard y cables de conexión.
    \end{enumerate}

    \section{Marco Teórico}

    \subsection{Osciloscopio}

    Un osciloscopio es un instrumento de visualización electrónico para la representación gráfica de señales eléctricas que pueden variar en el tiempo. Presenta los valores de las señales en forma de coordenadas en una pantalla, en la que normalmente el eje horizontal representa el tiempo y el eje vertical representa la tensión. La imagen así obtenida se denomina oscilograma~\cite{tektronix2009xyz}.

    A partir de una medición con el osciloscopio se puede obtener:

    \begin{itemize}
        \item El comportamiento en función del tiempo de las señales de tensión.
        \item Anomalías o distorsión en las señales de tensión.
        \item La frecuencia de una señal oscilante y cualquier variación en ella.
        \item Si la señal tiene ruido y los efectos que este produce sobre ella.
    \end{itemize}

    \subsection{Generador de funciones}

    Un generador de funciones es un generador de señal que permite producir múltiples señales eléctricas con formas y características ajustables, tales como la frecuencia, la amplitud y la forma de onda. Entre las formas de onda más comunes se encuentran:

    \begin{itemize}
        \item Senoidales.
        \item Cuadradas y rectangulares.
        \item Triangulares y diente de sierra.
        \item Escalón y pulso.
        \item Señales tipo TTL.
    \end{itemize}

    \subsection{Señales senoidales}

    Las señales senoidales son las más empleadas en circuitos eléctricos porque satisfacen las siguientes propiedades~\cite{alexander2018fundamentos}:

    \begin{itemize}
        \item Una red eléctrica lineal excitada por una fuente de tensión o intensidad senoidal muestra, en todas las partes de la red, tensiones e intensidades senoidales (pasado un corto período transitorio). Estas tienen la misma frecuencia del generador, pudiendo variar en fase y/o amplitud.
        \item La suma de $n$ señales senoidales de amplitudes y fases arbitrarias, pero de la misma frecuencia, es también una señal senoidal de igual frecuencia que la de las señales que la componen.
        \item La derivada de una señal senoidal es una señal senoidal.
        \item La integral de una señal senoidal es una señal senoidal.
    \end{itemize}

    \subsubsection{Parámetros que definen a la señal senoidal}

    Una señal senoidal se define por la siguiente ecuación:

    \begin{equation}
        X(t) = X_p \sin(\omega t + \theta)\label{eq:senoide}
    \end{equation}

    donde $X_p = A$ es la amplitud o valor pico, $\omega$ es la pulsación o frecuencia angular y $\theta$ es el ángulo de fase.

    \begin{figure}[ht!]
        \centering
        \begin{tikzpicture}
            \begin{axis}[
                    width=11cm, height=5.5cm,
                    axis lines=middle,
                    xlabel={$t$}, ylabel={$X(t)$},
                    xtick={0.5,1,1.5}, xticklabels={$T/2$,$T$,$3T/2$},
                    ytick={-1,1}, yticklabels={$-X_p$,$X_p$},
                    ymin=-1.5, ymax=1.5, xmin=0, xmax=1.65,
                    every axis x label/.style={at={(ticklabel* cs:1.02)}, anchor=west},
                    every axis y label/.style={at={(ticklabel* cs:1.02)}, anchor=south},
                    clip=false
                ]
                \addplot[thick, domain=0:1.6, samples=200] {sin(deg(2*pi*x))};
                \addplot[dashed, domain=0:1.6, samples=2] {0.7071};
                \node[anchor=south west, font=\small] at (axis cs:1.2,0.72) {$X_{RMS}$};
                \draw[<->] (axis cs:0.25,-1) -- (axis cs:0.25,1);
                \node[anchor=east, font=\small] at (axis cs:0.25,0.5) {$X_{pp}$};
            \end{axis}
        \end{tikzpicture}
        \caption{Señal senoidal con $\theta = 0$, valor pico $X_p$, valor pico a pico $X_{pp}$ y valor eficaz $X_{RMS}$}\label{fig:senoide}
    \end{figure}

    \subsubsection{Frecuencia angular}

    La frecuencia angular $\omega$ es el número de oscilaciones o vueltas por segundo que realiza la onda, medido en [rad/s]. Se relaciona con la frecuencia cíclica $f$ de la siguiente manera:

    \[
        \omega = 2\pi f
    \]

    \subsubsection{Período}

    El período $T$ de una función periódica es el tiempo de un ciclo completo. Su recíproco es el número de ciclos por segundo, conocido como frecuencia cíclica $f$:

    \[
        T = \frac{1}{f} = \frac{2\pi}{\omega}
    \]

    \subsubsection{Valor pico y valor pico a pico}

    El valor pico $X_p$ es el mayor valor instantáneo que toma la señal, también llamado magnitud o amplitud de la onda. El valor pico a pico $X_{pp}$ es la diferencia entre el valor máximo y el valor mínimo de la señal. Para una señal simétrica y sin componente continua:

    \[
        X_{pp} = 2X_p
    \]

    \subsubsection{Valor eficaz}

    El valor eficaz o valor cuadrático medio $X_{RMS}$ se define como el valor de una señal de corriente directa que produce la misma disipación de energía que su equivalente en corriente alterna en un período de la señal alterna. Matemáticamente, es la raíz cuadrada de la media de los cuadrados de los valores instantáneos alcanzados en un período:

    \begin{equation}
        X_{RMS} = \sqrt{\frac{1}{T}\int_{0}^{T} X(t)^2\,dt}\label{eq:rms}
    \end{equation}

    Para una señal \textbf{senoidal} la relación entre el valor pico y el valor eficaz es:

    \[
        V_p = \sqrt{2}\,V_{RMS}
    \]

    y para una señal \textbf{triangular}:

    \[
        V_p = \sqrt{3}\,V_{RMS}
    \]

    La demostración analítica de estas relaciones a partir de la ecuación~\eqref{eq:rms} se incluye en el Anexo~\ref{anexo:demostracion}.

    \subsubsection{Valor promedio}

    El valor promedio es un valor constante que representa la «media» de la onda, y para cualquier forma de onda está definido por:

    \[
        X_{prom} = \frac{1}{T}\int_{0}^{T} X(t)\,dt
    \]

    Para las señales senoidal y triangular simétricas respecto al eje horizontal, el valor promedio en un período es cero.

    \section{Procedimiento Experimental}

    \subsection{Montaje 1}

    \begin{enumerate}
        \item Montar el circuito de la figura~\ref{fig:montaje} utilizando el generador de funciones en lugar de una fuente de voltaje, con los datos de la tabla~\ref{tab:m1}.

        % >>> TODO (LLENAR): verifica en el generador si la amplitud se define como Vp o Vpp, y
        % que la carga esté configurada como "Hi-Z" (la resistencia de 75 kΩ es mucho mayor que
        % 50 Ω; con la carga mal configurada la amplitud real puede ser el doble de la mostrada).
        % Anota aquí la configuración usada si lo consideras relevante.

        \begin{figure}[ht!]
            \centering
            \begin{circuitikz}[american]
                \draw (0,0) to[sV, l=$V_S$] (0,3)
                      -- (3,3) to[R, l=$R$] (3,0)
                      -- (0,0) node[ground] {};
            \end{circuitikz}
            \caption{Circuito de los montajes 1 y 2}\label{fig:montaje}
        \end{figure}

        \begin{table}[ht!]
            \centering
            \begin{tabular}{ll}
                \toprule
                Parámetro      & Valor                      \\
                \midrule
                Voltaje pico   & \SI{7}{\volt}              \\
                Función        & Senoidal                   \\
                Frecuencia     & \SI{50}{\hertz}            \\
                Resistencia    & \SI{75}{\kilo\ohm}         \\
                \bottomrule
            \end{tabular}
            \caption{Datos del montaje 1}\label{tab:m1}
        \end{table}

        \item Preparar el osciloscopio y el multímetro para realizar mediciones de voltaje en la resistencia.

        % >>> TODO (LLENAR): si lo deseas, agrega una foto o captura del montaje real:
        % \imagen{./Pictures/Practica1/montaje1.jpg}{Montaje 1 implementado}{250px}

        \item Capturar la forma de onda mostrada en el osciloscopio, indicando la escala en V/div y s/div (figura~\ref{fig:cap-m1} en el Anexo~\ref{anexo:capturas}).

        \item Utilizando la lectura del osciloscopio, realizar las siguientes mediciones:

        % >>> TODO (LLENAR): mediciones experimentales del montaje 1 con el osciloscopio.
        \begin{table}[ht!]
            \centering
            \begin{tabular}{lc}
                \toprule
                Medición                         & Valor  \\
                \midrule
                Voltaje pico $V_P$ (V)           & \vacio \\
                Voltaje pico a pico $V_{PP}$ (V) & \vacio \\
                Período $T$ (ms)                 & \vacio \\
                \bottomrule
            \end{tabular}
            \caption{Mediciones con el osciloscopio, montaje 1}\label{tab:osc-m1}
        \end{table}

        \item Utilizando el multímetro, realizar la medición del voltaje eficaz en la resistencia:

        % >>> TODO (LLENAR): lectura del multímetro (en modo de voltaje AC) del montaje 1.
        \[
            V_{RMS} = \vacio\ \si{\volt}
        \]
    \end{enumerate}

    \subsection{Montaje 2}

    \begin{enumerate}
        \item Montar el circuito de la figura~\ref{fig:montaje} utilizando el generador de funciones en lugar de una fuente de voltaje, con los datos de la tabla~\ref{tab:m2}.

        \begin{table}[ht!]
            \centering
            \begin{tabular}{ll}
                \toprule
                Parámetro      & Valor                      \\
                \midrule
                Voltaje pico   & \SI{8}{\volt}              \\
                Función        & Triangular                 \\
                Frecuencia     & \SI{60}{\hertz}            \\
                Resistencia    & \SI{75}{\kilo\ohm}         \\
                \bottomrule
            \end{tabular}
            \caption{Datos del montaje 2}\label{tab:m2}
        \end{table}

        \item Preparar el osciloscopio y el multímetro para realizar mediciones de voltaje en la resistencia.

        \item Capturar la forma de onda mostrada en el osciloscopio, indicando la escala en V/div y s/div (figura~\ref{fig:cap-m2} en el Anexo~\ref{anexo:capturas}).

        \item Utilizando la lectura del osciloscopio, realizar las siguientes mediciones:

        % >>> TODO (LLENAR): mediciones experimentales del montaje 2 con el osciloscopio.
        \begin{table}[ht!]
            \centering
            \begin{tabular}{lc}
                \toprule
                Medición                         & Valor  \\
                \midrule
                Voltaje pico $V_P$ (V)           & \vacio \\
                Voltaje pico a pico $V_{PP}$ (V) & \vacio \\
                Período $T$ (ms)                 & \vacio \\
                \bottomrule
            \end{tabular}
            \caption{Mediciones con el osciloscopio, montaje 2}\label{tab:osc-m2}
        \end{table}

        \item Utilizando el multímetro, realizar la medición del voltaje eficaz en la resistencia:

        % >>> TODO (LLENAR): lectura del multímetro del montaje 2.
        % NOTA: si el multímetro NO es de verdadero valor eficaz (True RMS), su lectura
        % con una onda triangular no será confiable (la mayoría está calibrado para
        % senoides). Si pasa, anótalo y coméntalo como fuente de error en el análisis.
        \[
            V_{RMS} = \vacio\ \si{\volt}
        \]
    \end{enumerate}

    \subsection{Montaje 3}

    Se utilizan dos circuitos similares al de la figura~\ref{fig:montaje}, con el mismo valor de $R$.

    % >>> TODO (LLENAR): anota el valor de R utilizado en este montaje.
    El valor de resistencia utilizado en ambos circuitos fue $R = \vacio\ \si{\kilo\ohm}$.

    \begin{enumerate}
        \item Aplicar al primer circuito, con el generador de funciones, la señal del canal 1:

        \[
            V_1(t) = 2\sin(320\pi t)\ \si{\volt}
        \]

        \item Aplicar al segundo circuito, con el generador de funciones, la señal del canal 2:

        \[
            V_2(t) = 4\sin(320\pi t)\ \si{\volt}
        \]

        Como $\omega = 320\pi\ \si{\radian\per\second}$, ambas señales tienen una frecuencia $f = \omega/2\pi = \SI{160}{\hertz}$ y un período $T = 1/f = \SI{6.25}{\milli\second}$.

        \item Capturar la lectura del osciloscopio de forma simultánea de ambas señales, con el mismo osciloscopio (figura~\ref{fig:cap-m3} en el Anexo~\ref{anexo:capturas}).

        \item Con la función matemática adecuada del osciloscopio (CH1 + CH2), mostrar la señal resultante de sumar ambas señales (figura~\ref{fig:cap-m3-suma} en el Anexo~\ref{anexo:capturas}) y, a partir de ella, identificar:

        % >>> TODO (LLENAR): mediciones de la señal resultante (CH1 + CH2).
        \begin{table}[ht!]
            \centering
            \begin{tabular}{lc}
                \toprule
                Medición                         & Valor  \\
                \midrule
                Voltaje pico $V_P$ (V)           & \vacio \\
                Voltaje pico a pico $V_{PP}$ (V) & \vacio \\
                Período $T$ (ms)                 & \vacio \\
                \bottomrule
            \end{tabular}
            \caption{Mediciones de la señal resultante, montaje 3}\label{tab:osc-m3}
        \end{table}
    \end{enumerate}

    \section{Análisis de Resultados}

    Una vez terminada la práctica, se realizan los siguientes análisis. En todos los casos, el error relativo porcentual se calcula como:

    \[
        E_{\%} = \frac{\left|V_{exp} - V_{teor}\right|}{V_{teor}} \times 100\%
    \]

    \begin{enumerate}
        \item Comparación de los valores medidos y teóricos del \textbf{montaje 1}.

        \begin{quote}
            Como la fuente está conectada directamente a la resistencia, el voltaje en $R$ es el voltaje del generador, por lo que los valores teóricos se calculan a partir de los datos de la tabla~\ref{tab:m1}:

            \[
                V_P = \SI{7}{\volt}
            \]

            \[
                V_{PP} = 2V_P = 2(7) = \SI{14}{\volt}
            \]

            \[
                V_{RMS} = \frac{V_P}{\sqrt{2}} = \frac{7}{\sqrt{2}} \approx \SI{4.950}{\volt}
            \]

            \[
                T = \frac{1}{f} = \frac{1}{50} = \SI{20}{\milli\second}
            \]

            % >>> TODO (LLENAR): sustituir el valor experimental (primer espacio) en cada error relativo.
            Los errores relativos porcentuales son:

            \[
                E_{V_P} = \frac{\left|\,\vacio - 7\,\right|}{7} \times 100\% = \vacio\,\%
            \]

            \[
                E_{V_{PP}} = \frac{\left|\,\vacio - 14\,\right|}{14} \times 100\% = \vacio\,\%
            \]

            \[
                E_{V_{RMS}} = \frac{\left|\,\vacio - 4.950\,\right|}{4.950} \times 100\% = \vacio\,\%
            \]

            \[
                E_{T} = \frac{\left|\,\vacio - 20\,\right|}{20} \times 100\% = \vacio\,\%
            \]
        \end{quote}

        % >>> TODO (LLENAR): columnas "Valor experimental" y "Error" de la tabla.
        % NOTA: el manual rotula la fila del período como "T (Hz)", pero el período
        % se mide en segundos (ms); aquí se usa ms.
        \begin{table}[ht!]
            \centering
            \begin{tabular}{lccc}
                \toprule
                                & Valor experimental & Valor teórico & Error (\%)    \\
                \midrule
                $V_P$ (V)       & \vacio[1.8cm]      & 7             & \vacio[1.5cm] \\
                $V_{PP}$ (V)    & \vacio[1.8cm]      & 14            & \vacio[1.5cm] \\
                $V_{RMS}$ (V)   & \vacio[1.8cm]      & 4.950         & \vacio[1.5cm] \\
                $T$ (ms)        & \vacio[1.8cm]      & 20            & \vacio[1.5cm] \\
                \bottomrule
            \end{tabular}
            \caption{Resumen de resultados del montaje 1}\label{tab:res-m1}
        \end{table}

        \item Comparación de los valores medidos y teóricos del \textbf{montaje 2}.

        \begin{quote}
            Para la onda triangular se cumple $V_P = \sqrt{3}\,V_{RMS}$. Con los datos de la tabla~\ref{tab:m2}:

            \[
                V_P = \SI{8}{\volt}
            \]

            \[
                V_{PP} = 2V_P = 2(8) = \SI{16}{\volt}
            \]

            \[
                V_{RMS} = \frac{V_P}{\sqrt{3}} = \frac{8}{\sqrt{3}} \approx \SI{4.619}{\volt}
            \]

            \[
                T = \frac{1}{f} = \frac{1}{60} \approx \SI{16.667}{\milli\second}
            \]

            % >>> TODO (LLENAR): sustituir el valor experimental (primer espacio) en cada error relativo.
            Los errores relativos porcentuales son:

            \[
                E_{V_P} = \frac{\left|\,\vacio - 8\,\right|}{8} \times 100\% = \vacio\,\%
            \]

            \[
                E_{V_{PP}} = \frac{\left|\,\vacio - 16\,\right|}{16} \times 100\% = \vacio\,\%
            \]

            \[
                E_{V_{RMS}} = \frac{\left|\,\vacio - 4.619\,\right|}{4.619} \times 100\% = \vacio\,\%
            \]

            \[
                E_{T} = \frac{\left|\,\vacio - 16.667\,\right|}{16.667} \times 100\% = \vacio\,\%
            \]
        \end{quote}

        % >>> TODO (LLENAR): columnas "Valor experimental" y "Error" de la tabla.
        \begin{table}[ht!]
            \centering
            \begin{tabular}{lccc}
                \toprule
                                & Valor experimental & Valor teórico & Error (\%)    \\
                \midrule
                $V_P$ (V)       & \vacio[1.8cm]      & 8             & \vacio[1.5cm] \\
                $V_{PP}$ (V)    & \vacio[1.8cm]      & 16            & \vacio[1.5cm] \\
                $V_{RMS}$ (V)   & \vacio[1.8cm]      & 4.619         & \vacio[1.5cm] \\
                $T$ (ms)        & \vacio[1.8cm]      & 16.667        & \vacio[1.5cm] \\
                \bottomrule
            \end{tabular}
            \caption{Resumen de resultados del montaje 2}\label{tab:res-m2}
        \end{table}

        \item De acuerdo con las lecturas del \textbf{montaje 3}:

        \begin{enumerate}[label=\alph*)] % chktex 10
            \item ¿Cómo es la señal resultante con respecto de las señales descompuestas que se suman? ¿Es una señal senoidal?

            \begin{quote}
                % >>> TODO (LLENAR): verificar con la captura que lo observado coincide con esta respuesta.
                Sí, la señal resultante es senoidal. Como ambas señales tienen la misma frecuencia y la misma fase, su suma es otra senoide de igual frecuencia, cuya amplitud es la suma de las amplitudes:

                \[
                    V_1(t) + V_2(t) = 2\sin(320\pi t) + 4\sin(320\pi t) = 6\sin(320\pi t)\ \si{\volt}
                \]

                Es decir, la señal resultante está en fase con las originales y tiene una amplitud mayor, $V_P = \SI{6}{\volt}$. Esto concuerda con la propiedad de las señales senoidales mencionada en el marco teórico.
            \end{quote}

            \item De la señal resultante, ¿es posible identificar voltaje pico a pico, frecuencia y período a partir de la lectura del osciloscopio? ¿Por qué?

            \begin{quote}
                Sí, es posible. Al ser una señal senoidal periódica con una forma bien definida, el voltaje pico a pico se lee como la distancia vertical entre el máximo y el mínimo (multiplicada por los V/div), el período como la distancia horizontal de un ciclo completo (multiplicada por los s/div), y la frecuencia se obtiene como $f = 1/T$.

                Teóricamente se espera $V_{PP} = 2(6) = \SI{12}{\volt}$ y $T = \SI{6.25}{\milli\second}$.
            \end{quote}

            \item Calcular la frecuencia de la señal resultante a partir de las mediciones obtenidas y comparar los valores de frecuencia de las tres señales (canal 1, canal 2 y resultante).

            \begin{quote}
                La frecuencia teórica de las tres señales es:

                \[
                    f = \frac{\omega}{2\pi} = \frac{320\pi}{2\pi} = \SI{160}{\hertz}
                \]

                % >>> TODO (LLENAR): período medido de la señal resultante y frecuencia calculada.
                A partir del período medido en la señal resultante:

                \[
                    f = \frac{1}{T} = \frac{1}{\vacio\ \si{\second}} = \vacio\ \si{\hertz}
                \]

                % >>> TODO (LLENAR): período medido de cada señal, su frecuencia y el error respecto a 160 Hz.
                \begin{table}[ht!]
                    \centering
                    \begin{tabular}{lccc}
                        \toprule
                        Señal         & $T$ (ms)      & $f = 1/T$ (Hz) & Error (\%)    \\
                        \midrule
                        Canal 1       & \vacio[1.5cm] & \vacio[1.5cm]  & \vacio[1.5cm] \\
                        Canal 2       & \vacio[1.5cm] & \vacio[1.5cm]  & \vacio[1.5cm] \\
                        Resultante    & \vacio[1.5cm] & \vacio[1.5cm]  & \vacio[1.5cm] \\
                        \bottomrule
                    \end{tabular}
                    \caption{Frecuencias medidas en el montaje 3 (valor teórico: \SI{160}{\hertz})}\label{tab:freq-m3}
                \end{table}

                % >>> TODO (LLENAR): comentar si las tres frecuencias coinciden entre sí (dentro del error de medición).
                Se espera que las tres frecuencias coincidan, ya que al sumar señales senoidales de la misma frecuencia no se modifica la frecuencia, únicamente la amplitud y la fase.
            \end{quote}

            \item Si las señales de entrada tuvieran distintas frecuencias, ¿qué forma presentaría la señal resultante? ¿Sería posible identificar voltaje pico a pico, frecuencia y período a partir de la lectura del osciloscopio? ¿Por qué?

            \begin{quote}
                La señal resultante ya no sería una senoide, sino una forma de onda más compleja, producto de la superposición de ambas componentes. Será periódica únicamente si el cociente entre las frecuencias es un número racional; en ese caso su período es el mínimo tiempo en el que ambas componentes completan un número entero de ciclos, y su frecuencia es la frecuencia fundamental (el máximo común divisor de las dos frecuencias). Si el cociente es irracional, la señal no es periódica.

                Aun así, sí es posible leer el voltaje pico a pico (diferencia entre el máximo y el mínimo observados, que como mucho es la suma de los picos de ambas señales) y, cuando la señal es periódica, el período a partir de la repetición de la forma. Sin embargo, esos valores ya no corresponden a una única frecuencia ni a una amplitud senoidal, y su identificación es más difícil, sobre todo cuando las frecuencias son muy distintas o su cociente no es simple.
            \end{quote}
        \end{enumerate}

        \item Investigar y detallar las aplicaciones de cada tipo de onda disponible en el generador de funciones.

        \begin{quote}
            % >>> TODO (LLENAR): revisar y ampliar con las fuentes que realmente consultaste.
            % Verifica que las entradas de references.bib (Horowitz y Hill, Tektronix) coincidan
            % con lo que usaste, o reemplázalas por tus fuentes.
            \begin{itemize}
                \item \textbf{Senoidal:} es la forma de onda de la red eléctrica de corriente alterna y la base del análisis de circuitos en régimen permanente. Se usa para caracterizar la respuesta en frecuencia de circuitos y amplificadores, como señal de prueba en audio y como portadora en comunicaciones~\cite{alexander2018fundamentos,horowitz2015art}.
                \item \textbf{Cuadrada y rectangular:} se usa como señal de reloj en circuitos digitales, para estudiar la respuesta transitoria y el ancho de banda de circuitos (una onda cuadrada contiene armónicos de alta frecuencia), para compensar las puntas de prueba del osciloscopio y como señal de control tipo PWM en fuentes conmutadas~\cite{horowitz2015art,tektronix2009xyz}.
                \item \textbf{Triangular y diente de sierra:} se usa como señal de barrido (la base de tiempo de un osciloscopio es una rampa), para generar señales PWM al compararlas con una referencia, en osciladores controlados por voltaje, en pruebas de linealidad y en síntesis de sonido~\cite{horowitz2015art,tektronix2009xyz}.
                \item \textbf{Escalón y pulso:} se usan para estudiar la respuesta transitoria de circuitos (por ejemplo, la carga y descarga de circuitos RC, RL y RLC), como señales de disparo (trigger) y de temporización, y para probar sistemas de control~\cite{alexander2018fundamentos,horowitz2015art}.
                \item \textbf{TTL:} es una señal cuadrada con niveles lógicos compatibles con la familia TTL (aproximadamente \SI{0}{\volt} y \SI{5}{\volt}), utilizada para excitar y probar circuitos digitales, contadores, microcontroladores y como reloj de sincronización~\cite{horowitz2015art}.
            \end{itemize}
        \end{quote}

        \item Explicar la diferencia entre voltaje eficaz, voltaje pico y voltaje pico a pico.

        \begin{quote}
            % >>> TODO (LLENAR): revisar y complementar con las fuentes consultadas.
            El \textbf{voltaje pico} ($V_P$) es el mayor valor instantáneo que alcanza la señal medido respecto a su nivel de referencia (cero, en una señal sin componente continua). El \textbf{voltaje pico a pico} ($V_{PP}$) es la diferencia entre el valor máximo y el valor mínimo de la señal; para una señal simétrica, $V_{PP} = 2V_P$. Ambos son valores instantáneos extremos y se leen directamente en la pantalla del osciloscopio.

            El \textbf{voltaje eficaz} ($V_{RMS}$), en cambio, es un valor «promedio» en términos de energía: equivale al voltaje de corriente directa que disiparía la misma potencia en una resistencia~\cite{alexander2018fundamentos}. Depende de la forma de onda completa y no solo de su máximo: para una senoide $V_{RMS} = V_P/\sqrt{2} \approx 0.707\,V_P$ y para una onda triangular $V_{RMS} = V_P/\sqrt{3} \approx 0.577\,V_P$. Por ejemplo, la red de \SI{120}{\volt} eficaces tiene un valor pico de aproximadamente \SI{170}{\volt}. El multímetro suele indicar el valor eficaz, mientras que el osciloscopio permite leer los valores pico y pico a pico.
        \end{quote}
    \end{enumerate}

    \section{Conclusiones}

    % >>> TODO (LLENAR): redactar las conclusiones con base en los resultados obtenidos. Los
    % siguientes puntos son una guía; ajusta las frases y completa los espacios en blanco
    % según tus errores porcentuales y observaciones reales.
    \begin{enumerate}
        \item Se comprobó que el osciloscopio permite identificar el valor pico, el valor pico a pico y el período de una señal, mientras que el multímetro entrega el valor eficaz. En el montaje 1 el error relativo del voltaje eficaz fue de \vacio\,\%, por lo que \vacio\ (concluir si el valor eficaz experimental concuerda con la relación $V_p = \sqrt{2}\,V_{RMS}$).

        \item En el montaje 2, con una señal triangular, el error en el voltaje eficaz fue de \vacio\,\%, lo que \vacio\ (concluir si se verificó la relación $V_p = \sqrt{3}\,V_{RMS}$ y comentar posibles fuentes de error, como el tipo de multímetro utilizado).

        \item En el montaje 3 se observó que la suma de dos señales senoidales de la misma frecuencia es otra señal senoidal de igual frecuencia, con una frecuencia medida de \vacio\ \si{\hertz} frente a los \SI{160}{\hertz} esperados.

        \item \vacio\ (conclusión sobre el manejo del generador de funciones y del osciloscopio, y las dificultades o fuentes de error encontradas en la práctica).
    \end{enumerate}

    \nocite{unah2026lie311p1}
    \printbibliography{}

    \newpage
    \section{Anexos}

    \subsection{Demostración analítica de $V_p$ y $V_{RMS}$}\label{anexo:demostracion}

    \subsubsection{Señal senoidal}

    Sea $X(t) = X_p\sin(\omega t)$ con $\theta = 0$ y un período constante arbitrario $T = 2\pi/\omega$. A partir de la ecuación~\eqref{eq:rms}:

    \[
        X_{RMS}^2 = \frac{1}{T}\int_{0}^{T} X_p^2\sin^2(\omega t)\,dt
        = \frac{X_p^2}{T}\int_{0}^{T} \frac{1 - \cos(2\omega t)}{2}\,dt
    \]

    donde se usó la identidad $\sin^2\alpha = \dfrac{1 - \cos 2\alpha}{2}$. Integrando:

    \[
        X_{RMS}^2 = \frac{X_p^2}{2T}\left[t - \frac{\sin(2\omega t)}{2\omega}\right]_{0}^{T}
        = \frac{X_p^2}{2T}\left[T - \frac{\sin(2\omega T)}{2\omega}\right]
    \]

    Como $\omega T = 2\pi$, se tiene $\sin(2\omega T) = \sin(4\pi) = 0$, y por tanto:

    \[
        X_{RMS}^2 = \frac{X_p^2}{2T}\,T = \frac{X_p^2}{2}
        \quad\Longrightarrow\quad
        X_{RMS} = \frac{X_p}{\sqrt{2}}
        \quad\Longrightarrow\quad
        \boxed{X_p = \sqrt{2}\,X_{RMS}}
    \]

    El resultado es independiente del valor de $T$ elegido.

    \subsubsection{Señal triangular}

    Una onda triangular simétrica de valor pico $X_p$ y período $T$ está formada por tramos lineales de duración $T/4$ que van de $0$ a $X_p$ (y análogamente en los demás cuartos de período, cambiando solo el signo o el sentido). Por simetría, el valor cuadrático medio en un período es igual al de un solo tramo. En el primer cuarto, $X(t) = \dfrac{4X_p}{T}\,t$ para $0 \le t \le T/4$:

    \[
        X_{RMS}^2 = \frac{1}{T/4}\int_{0}^{T/4}\left(\frac{4X_p}{T}\,t\right)^2 dt
        = \frac{4}{T}\cdot\frac{16X_p^2}{T^2}\cdot\frac{(T/4)^3}{3}
        = \frac{X_p^2}{3}
    \]

    \[
        X_{RMS} = \frac{X_p}{\sqrt{3}}
        \quad\Longrightarrow\quad
        \boxed{X_p = \sqrt{3}\,X_{RMS}}
    \]

    \subsection{Capturas del osciloscopio}\label{anexo:capturas}

    % >>> TODO (LLENAR): reemplazar cada \capturapendiente por la captura real, por ejemplo:
    % \imagen{./Pictures/Practica1/cap_m1.png}{Captura del montaje 1 (V/div = ..., s/div = ...)}{475px}
    % Recuerda indicar siempre la escala (V/div y s/div) en cada captura.
    % Si reemplazas el recuadro, las referencias \ref{fig:cap-m1}, etc. dejarán de funcionar
    % a menos que agregues \label{fig:cap-m1} (y los demás) al pie de la nueva figura.

    \capturapendiente{Montaje 1: señal senoidal en $R$}{V/div: \vacio[1.2cm] \quad s/div: \vacio[1.2cm]}{fig:cap-m1}

    \capturapendiente{Montaje 2: señal triangular en $R$}{V/div: \vacio[1.2cm] \quad s/div: \vacio[1.2cm]}{fig:cap-m2}

    \capturapendiente{Montaje 3: canal 1 y canal 2 de forma simultánea}{V/div: \vacio[1.2cm] \quad s/div: \vacio[1.2cm]}{fig:cap-m3}

    \capturapendiente{Montaje 3: señal resultante de la suma (CH1 + CH2)}{V/div: \vacio[1.2cm] \quad s/div: \vacio[1.2cm]}{fig:cap-m3-suma}

\end{document}